\documentclass[11pt]{article}

\usepackage{graphicx}
\usepackage{xcolor}
\usepackage{geometry}
\usepackage{amsmath}
\usepackage{amssymb}
\usepackage{booktabs}
\usepackage{array}
\usepackage{enumitem}
\usepackage{algorithm}
\usepackage{algpseudocode}

\geometry{letterpaper, top=2.5cm, bottom=2.5cm, left=3cm, right=2.5cm}
\emergencystretch=3em
\setlist[itemize]{itemsep=0.3em, leftmargin=1.5em}
\setlist[enumerate]{itemsep=0.3em, leftmargin=1.8em}

\usepackage[
    colorlinks=true,
    linkcolor=blue,
    citecolor=darkgray,
    urlcolor=blue,
    bookmarks=true,
    bookmarksnumbered=true,
    unicode=true
]{hyperref}
\hypersetup{
    pdftitle={PULSE with SGD: Implementation Notes (pool + batch version)},
    pdfsubject={PULSE peptide-observation model, stochastic gradient descent},
    pdfkeywords={PULSE, SGD, proposal distribution, emergence support}
}

\title{\textbf{PULSE with SGD: Implementation Notes\\ (pool + batch version)}}
\date{\today}

\begin{document}
\maketitle

This document describes how the PULSE pipeline (\texttt{pulse/})
implements stochastic gradient descent for the PULSE peptide-observation
model. It follows the structure and notation of the June reference writeup
(\texttt{SGD\_PULSE\_reference.pdf}, the one-batch version) and highlights
where the current pool + batch implementation differs.

The two structural changes relative to June are:
\begin{enumerate}
\item \textbf{Pool + batch surveillance design.} The observed nucleotide
table is split by time into a \textbf{pool} (the first whole time columns
holding $\sim$\texttt{pool.fraction} of all counts, default 1\%) and a
\textbf{batch} (everything later, plus the \texttt{unknown} column).
Emergence support always flows from pool parents; the batch holds the
sequences whose emergence we explain.
\item \textbf{Multi-mutation proposal.} The unobserved space is sampled
with up to \texttt{n\_max} (default 3) simultaneous amino-acid mutations
per draw, with the mutation count drawn from a truncated Poisson. The
proposal probability $q$ and its rejection normalization $Z$ are computed
exactly.
\end{enumerate}

\section{Likelihood}\label{sec:likelihood}

Let $\mathcal{D}^Y_{\text{pool}}$ be the set of pool nucleotide sequences
(parents), $\mathcal{D}^Y_{\text{obs}}$ the set of observed batch
nucleotide sequences that pass the distance filter
(Section~\ref{sec:observed}), and $\mathcal{D}^X_{\text{obs}}$ the
corresponding observed amino-acid sequences. For an amino-acid sequence
$x$, let $\hat{\mathcal{Y}}(x)$ be its tractable encoding set
(Section~\ref{sec:yhat}). The approximated log-likelihood is
\begin{equation}\label{eq:likelihood}
\hat{\ell}(\theta, \beta, p_o)
= \sum_{x \in \mathcal{D}^X_{\text{obs}}}
\log \Big( p_a(x; \theta) \, s(x; \beta, p_o) \Big)
+ \sum_{x' \in \mathcal{U}}
\log \Big( 1 - p_a(x'; \theta) \, s(x'; \beta, p_o) \Big),
\end{equation}
where the \textbf{detection probability} of any peptide $x$ is
\begin{equation}\label{eq:detection}
s(x; \beta, p_o)
= 1 - \prod_{y \in \hat{\mathcal{Y}}(x)}
\big(1 - p_o \, p_e(y; \beta)\big),
\qquad
p_e(y; \beta) = 1 - \exp\big(-\beta \, S(y)\big),
\end{equation}
$p_a(x;\theta)$ is the activity probability from the classifier (LR or
NAM) over peptide features, $p_o$ the observation probability, $S(y)$ the
emergence support of nucleotide $y$ from the pool
(Section~\ref{sec:support}), and $\mathcal{U}$ the unobserved peptide
space --- every peptide reachable from the pool that is \textbf{not} any
table peptide.

Correspondence with the June notation: $\beta$ plays the role of
$\alpha T$ (the population scale $T$ is absorbed into $\beta$), and pool
frequency $\text{freq}_0(y)$ plays the role of $f(y)$.

The unobserved sum over $\mathcal{U}$ cannot be enumerated, so it is
estimated by importance sampling (Section~\ref{sec:estimator}). The
observed sum is computed exactly at every step.

\section{Why the unobserved sum must be importance-sampled}\label{sec:why}

The June writeup (its Section 3) showed that naively subsampling observed
sequences and mutating only those does \textbf{not} give an unbiased
gradient, because the same unobserved $x'$ would carry different encoding
sets $\hat{\mathcal{Y}}(x')$ in different minibatches. The fix is the same
here: for each \textbf{sampled} $x'$ we compute its contribution
\textbf{deterministically} --- its encoding set, the support of each
encoding, and its proposal probability $q(x')$ --- from one-time pool
structures that never change during training. The estimator
(Section~\ref{sec:estimator}) is then unbiased for the finite sum over
$\mathcal{U}$.

\section{Observed side (one-time precompute, cached)}\label{sec:observed}

All of Section~\ref{sec:observed} is computed once per configuration
--- data, pool split, rate model, and the two distance knobs
\texttt{obs\_dist\_max} and \texttt{support\_max\_dist} --- and cached to
disk. The cached pieces are \texttt{OBS\_L0}, \texttt{OBS\_L2}, and
\texttt{POS\_L1} files under \texttt{support\_cache/}; training reruns
load them.

\subsection{Pool / batch split}\label{sec:split}

Time columns are recognized as \texttt{YYYY}, \texttt{YYYY-MM}, or
\texttt{YYYY-MM-DD} (zero-padded, so string order $=$ chronological
order), plus a literal \texttt{before 2000} column that sorts first.
Walking columns in chronological order, the pool takes whole columns up to
and including the first column whose cumulative count reaches
\texttt{pool.fraction} of the grand total (whole columns are kept even
when they overshoot the fraction). Every later column plus
\texttt{unknown} is batch. Each row then gets \texttt{Pool0\_Count} and
\texttt{Batch\_Count}.

Row filters applied at load time, each reporting its drop count:
\begin{itemize}
\item \textbf{length filter}: rows whose nucleotide or peptide length
differs from the first row's are dropped;
\item \textbf{zero-count filter}: rows with no counts in any time column
or \texttt{unknown} are dropped;
\item \textbf{stop codons}: nucleotide sequences translating with an
internal \texttt{*} are excluded from supports (they never become parents
or encodings);
\item optional \texttt{drop\_gap\_or\_ambiguous} (off by default): rows
whose nucleotide string is not pure \texttt{ACGT} are dropped. By default,
\texttt{-{}-{}-}/\texttt{NNN} rows are kept: gap/ambiguous sites are simply
excluded from distance and rate computations on both sides (``mutually
valid sites''), and the encoded amino acids \texttt{X}/\texttt{-} carry a
$(0,0,0)$ descriptor row.
\end{itemize}

\subsection{Support of an observed batch sequence}\label{sec:support}

A unique batch nucleotide sequence $y$ (Batch\_Count $> 0$) is
\textbf{kept} if its minimal distance to the pool is $\le$
\texttt{obs\_dist\_max} (default 2; distance $=$ Hamming over mutually
valid sites) and \textbf{excluded} otherwise --- exclusions are reported.

The support of a kept sequence aggregates every independent way the pool
can produce it:
\begin{equation}\label{eq:support}
S(y) \;=\;
\underbrace{\sum_{\substack{y_p \in \mathcal{D}^Y_{\text{pool}}\\
1 \le d(y_p, y) \le d_S}}
P_{\min}(y_p \!\to\! y)\; \text{freq}_0(y_p)}_{\text{mutation from other parents}}
\;+\;
\underbrace{\mathbf{1}[y \in \text{pool}]\;
P_{\text{self}}(y)\; \text{freq}_0(y)}_{\text{self-survival (exact match only)}}
\end{equation}
with $d_S$ = \texttt{support\_max\_dist} (default 3) and
\begin{equation}\label{eq:pmin}
P_{\min}(y_p \!\to\! y) = \prod_{i:\, y_p[i] \ne y[i],\; \text{both valid}}
\mu\big(y_p[i] \!\to\! y[i]\big),
\qquad
P_{\text{self}}(y) = \prod_{i \text{ valid}}
\Big(1 - \sum_{b' \ne y[i]} \mu\big(y[i] \!\to\! b'\big)\Big).
\end{equation}
$P_{\min}$ is the minimal-pathway rate product: because per-site mutation
rates are tiny, paths with extra mutations are negligible, so the direct
product over the differing sites dominates and is used exactly. Note the
two distances are independent knobs: \texttt{obs\_dist\_max} decides
\textbf{which} batch sequences are observed; \texttt{support\_max\_dist}
decides \textbf{which parents} may contribute support (kept sequences
farther than \texttt{support\_max\_dist} from all parents get mutation
support 0 --- only self-survival if applicable). A distance-0 ``parent''
that is merely a gap/ambiguity variant of $y$ (same over mutually valid
sites, different string) contributes nothing --- matching the rule that
self-support is exact-match only.

Per-site rates $\mu(b \!\to\! b')$ come from the mutation model:
\begin{itemize}
\item default \texttt{ti90\_tv5}, \emph{per-base} normalization:
transition $1.8\times10^{-4}$, transversion $1\times10^{-5}$, so
ti $+$ 2\,tv $= 2\times10^{-4}$ at every site;
\item opt-in \texttt{sars\_cov2}, \emph{spectrum} normalization: the
12-pathway weight table (C$\to$T 46.5, G$\to$T 18.2, A$\to$G 9.4,
T$\to$C 9.4, G$\to$A 7.4, the remaining seven transversions 1.3) scaled
so the genome-average per-nucleotide rate is $\mu_{\text{total}} =
2\times10^{-4}$ (rate $= 4\mu_{\text{total}} w / \sum_{12} w$).
Cross-pathway ratios then equal the weight ratios exactly; per-base
totals differ (C sites mutate fastest).
\end{itemize}

\subsection{Peptide aggregation and $\hat{\mathcal{Y}}(x)$}\label{sec:yhat}

Observed batch nucleotides are translated and grouped; \textbf{peptide
support is the sum over its observed encodings} (each encoding's $S(y)$
already bundles its own self-survival term).

For the detection product, $\hat{\mathcal{Y}}(x)$ of an observed peptide
spans not only its observed batch encodings but also \textbf{extra
encodings} (computed once, cached as \texttt{POS\_L1}):
\begin{itemize}
\item \textbf{enc2}: minimal-codon-path encodings of $x$ from pool
parents whose peptide is at amino-acid distance $1 \dots n_{\max}$ from
$x$ (same machinery and nucleotide budget as the sampler,
Section~\ref{sec:encodings});
\item \textbf{enc3}: when $x$ is itself a pool peptide, single-position
synonymous codon variants ($k = 1 \dots n_{\max}$ changes in one codon)
of its pool member nucleotides.
\end{itemize}
Table nucleotide strings are excluded from the extras (observed encodings
already enter with their own support). Every extra encoding carries the
exact support of Section~\ref{sec:support}.

\section{The proposal distribution over unobserved peptides}\label{sec:proposal}

There is no candidate enumeration and no support threshold anywhere on
the unobserved side: the space is sampled, never materialized.

\subsection{One raw draw}\label{sec:draw}
\begin{enumerate}
\item \textbf{Root}: sample a pool peptide $x \sim r$, where $r$ is the
pool-frequency distribution aggregated over synonymous pool nucleotides:
\begin{equation}
r(x) = \sum_{\substack{y_p \in \text{pool}\\ \text{pep}(y_p)=x}}
\text{freq}_0(y_p).
\end{equation}
Roots with no canonical residue (all \texttt{X}/\texttt{-}/\texttt{*})
get weight 0.
\item \textbf{Mutation count}: $n \sim \text{Poisson}(\lambda)$
truncated to $\{1, \dots, \min(n_{\max}, L_x)\}$ and renormalized per
root, where $L_x$ is the number of canonical residues of $x$ and
$\lambda = L_{\text{nuc}} \cdot \mu_{\text{total}}$ by default
(\texttt{proposal.lam: auto}; any positive float overrides). This is the
expected number of nucleotide mutations per sequence per replication
cycle, so the draw mass is concentrated at $n = 1$ with a tail decaying
at scale $\mu_{\text{total}}$.
\item \textbf{Mutation}: pick $n$ distinct canonical positions of $x$
uniformly without replacement; at each, the new residue is uniform among
the other 19.
\item \textbf{Rejection}: if $x'$ is \textbf{any table peptide} (a
translation of any nucleotide row of the training file --- pool, batch,
or \texttt{unknown}), reject and resample.
\end{enumerate}

\subsection{Exact proposal probability}\label{sec:q}

For a parent peptide $p$ at amino-acid distance
$m \in \{1, \dots, n_{\max}\}$ from $x'$, the probability that one raw
draw from $p$ produces $x'$ is
\begin{equation}\label{eq:qmut}
q_{\text{mut}}(x' \mid p) \;=\; \frac{P_p(n = m)}
{\binom{L_p}{m}\, 19^{m}},
\qquad
P_p(n=m) = \frac{\lambda^m / m!}
{\sum_{j=1}^{\min(n_{\max}, L_p)} \lambda^j / j!},
\end{equation}
and the raw proposal mass of $x'$ sums over \textbf{all} candidate
parents --- the actual root plus every pool peptide at distance
$1 \dots n_{\max}$:
\begin{equation}\label{eq:qraw}
q_{\text{raw}}(x') \;=\; \sum_{m=1}^{n_{\max}}
\sum_{p:\, d_{aa}(p, x') = m} r(p)\, q_{\text{mut}}(x' \mid p).
\end{equation}
Omitting any parent would make the importance weight wrong and the
estimator biased --- this is the same requirement as in the June writeup,
extended to $m > 1$. (Implementation: candidate parents come from a
one-time pool$\times$pool amino-acid neighbor table with cutoff
$2 n_{\max}$ --- a draw from root $x$ with $n \le n_{\max}$ mutations can
only coincide with a parent within $2 n_{\max}$ of $x$.)

After rejection, the normalized proposal is
\begin{equation}\label{eq:z}
q(x') = \frac{q_{\text{raw}}(x')}{Z},
\qquad
Z = 1 - q_{\text{raw}}(\text{table})
= 1 - \sum_{p,\, m}
\frac{r(p)\, P_p(n=m)}{\binom{L_p}{m} 19^{m}} \, N_{\text{table}}(p, m),
\end{equation}
where $N_{\text{table}}(p, m)$ $=$ number of table peptides at amino-acid
distance $m$ from pool peptide $p$. $Z$ is computed \textbf{exactly} from
a one-time pool-peptide $\times$ table-peptide distance pass. Only when
that product is too large ($> 3\times10^8$ pairs) does a Monte-Carlo
pilot ($K = 200{,}000$ raw draws) estimate $Z$ instead; the mode is
printed at startup. A warning is printed if $Z < 0.5$ (more than half of
all draws rejected).

\subsection{Encodings of a sampled peptide and their support}\label{sec:encodings}

For an accepted draw $x'$ we deterministically build its encoding set
$\hat{\mathcal{Y}}(x')$:
\begin{itemize}
\item \textbf{Parents}: the actual root plus every pool peptide at
amino-acid distance $1 \dots n_{\max}$ from $x'$.
\item \textbf{Codon paths}: for parent peptide $p$, member nucleotide
$y_p$, and differing-residue positions $J$: mutations are confined to
the codons of $J$; at position $j$ any target codon of the new residue
is allowed; the total number of nucleotide changes over all of $J$ must
be $\le n_{\max}$. If $y_p$'s actual codon at $j$ cannot reach the
target residue within budget, that parent contributes \textbf{zero}
encodings --- codon context is respected exactly.
\item \textbf{Support}: each unique encoding $y'$ receives the exact
multi-distance support of Section~\ref{sec:support} --- the sum over all
pool parents at nucleotide distance $1 \dots$ \texttt{support\_max\_dist}
of $P_{\min} \cdot \text{freq}_0$ (candidate parents are found by a
triangle-inequality neighbor search around the generating parent at
cutoff $n_{\max} + d_S$). The same $y'$ reached from two parents is
merged once into the set; its support already aggregates all parents, so
nothing is double-counted.
\end{itemize}

A draw with no reachable encoding has $s(x') = 0$ \textbf{exactly} and
contributes $g = \log(1 - 0) = 0$: a wasted draw, but the estimator
stays unbiased.

\section{The SGD estimator}\label{sec:estimator}

Writing the unobserved contribution of peptide $x'$ as
\begin{equation}
g(x'; \theta, \beta, p_o)
= \log\big(1 - p_a(x'; \theta)\, s(x'; \beta, p_o)\big),
\end{equation}
the unobserved part of the likelihood is the finite sum
$\hat{L}_{\text{unobs}} = \sum_{x' \in \mathcal{U}} g(x')$.
With $B$ accepted draws $x'_b \sim q$ per step,
\begin{equation}\label{eq:estimator}
\widetilde{L}_{\text{unobs}} = \frac{1}{B} \sum_{b=1}^{B}
\frac{g(x'_b)}{q(x'_b)},
\qquad
\mathbb{E}_q\big[\widetilde{L}_{\text{unobs}}\big]
= \hat{L}_{\text{unobs}},
\end{equation}
because $q(x') > 0$ for every $x' \in \mathcal{U}$ (any reachable peptide
has positive raw mass and $Z > 0$) and $q(x') = 0$ on table peptides by
rejection. Since each $g(x'_b)$ is computed deterministically from the
one-time structures, the stochastic gradient is unbiased for the
approximated finite-sum objective.

The full per-step objective minimized by the optimizer is
\begin{equation}\label{eq:objective}
-\Bigg[
\underbrace{\sum_{x \in \mathcal{D}^X_{\text{obs}}}
\big(\log p_a(x) + \log s(x)\big)}_{\text{exact, every step}}
\;+\;
\underbrace{\frac{1}{B}\sum_{b=1}^{B}
\frac{\log\big(1 - p_a(x'_b)\, s(x'_b)\big)}{q(x'_b)}}_{\text{importance-sampled}}
\Bigg]
\;+\;
\lambda_{\text{reg}}\Big(\sum_{w \in \theta} w^2 + p_o^2\Big).
\end{equation}
One term per observed peptide, regardless of how many times it was seen
in surveillance (surveillance counts entered earlier, through the pool
frequencies and the pool/batch split --- this is the September one-batch
convention, deliberately different from June).

\section{Pseudocode}\label{sec:pseudocode}

Algorithm~\ref{alg:sgd} gives the full procedure: the one-time cached
precomputes (steps 1--6), then the per-step training loop (steps 7--24).

\begin{algorithm}[htbp]
\caption{PULSE SGD (pool + batch, multi-mutation proposal)}
\label{alg:sgd}
\begin{algorithmic}[1]
\Require training counts table, benchmark table, config
\Statex \textbf{One-time (cached):}
\State filter rows (length, zero-count, stop codons); split pool / batch
\State pool parents: unique pool nucleotides, $\text{freq}_0$, codon
table, peptide grouping
\State \texttt{OBS\_L0}: table $\times$ pool nucleotide neighbor pass
(cutoff $\max$ of the two knobs \texttt{obs\_dist\_max} and
\texttt{support\_max\_dist}); true minimal distance per table row
\State \texttt{OBS\_L2}: keep batch sequences with min distance $\le$
\texttt{obs\_dist\_max}; support $S(y)$ per kept sequence (mutation sum
$+$ self-survival); group by peptide; peptide support $=$ sum over
encodings
\State \texttt{SAMP\_L1}: root weights $r(x)$; pool--pool amino-acid
neighbor CSR (cutoff $2n_{\max}$); pool--pool nucleotide neighbor CSR
(cutoff $n_{\max} + d_S$); table-peptide set; $N_{\text{table}}$ counts
$\to$ exact $Z$ (or MC pilot)
\State \texttt{POS\_L1}: extra encodings (enc2, enc3) of observed
peptides
\Statex \textbf{Training:}
\State draw the \textsc{fixed} validation batch once
(\texttt{val\_seed}, 500 draws)
\For{step $t = 1, 2, \dots$}
\State batch $\gets B$ accepted draws (default $B = 128$):
\Repeat
\State $x \sim r$ ; \; $n \sim \text{truncPoisson}(\lambda \mid x)$
\State $n$ distinct canonical positions; new residue $\sim$
uniform(19)
\Until{$x'$ not a table peptide}
\State $q(x') = \big[ W(x, n) + \sum_{\text{AA-neighbor parents}}
W(p, m) \big] / Z$
\State $\hat{\mathcal{Y}}(x') =$ minimal-codon-path encodings (budget
$n_{\max}$); exact support per encoding
\State loss $= -(\, \text{positive\_term} + \text{mean}_b[\,
\log(1 - p_a s) / q \,]\,) + \text{reg}$
\State AdamW step on $(\theta, d_{raw}, b_{raw})$; grad-norm clip 1.0
\If{$t \bmod 10 = 0$}
\State validate on the fixed batch; keep best checkpoint
\State stop after \texttt{early\_stop\_checks} (default 100) checks
without \texttt{min\_delta} ($10^{-4}$) improvement
\EndIf
\EndFor
\State restore best checkpoint; report $\beta = \exp(d_{raw})$,
$p_o = \text{sigmoid}(b_{raw})$
\State score benchmark with $p_a$; AUC / AP / Spearman; append results
row; save model \texttt{.pt} $+$ predictions \texttt{.csv} under
\texttt{models/<run\_tag>}
\end{algorithmic}
\end{algorithm}

\section{Parameterization and numerical guards}\label{sec:guards}

\begin{itemize}
\item $\beta = \exp(d_{raw})$ --- free (no bound), initialized at
$d_{raw} = -\log(\text{median of positive supports})$ so that $\beta S$
starts in a gradient-sensitive zone.
\item $p_o = \text{sigmoid}(b_{raw})$, initialized at 0.5.
\item Probability clamps: $p_e, p_o, p_a \in [\texttt{p\_floor}, 1]$
with \texttt{p\_floor} $= 10^{-8}$ by default; $\log(1 - p_o p_e)$ is
clamped at \texttt{p\_floor}. These clamps are applied
\textbf{uniformly on both sides of the likelihood} --- observed peptides
and sampled unobserved draws flow through the same detection
computation, so every encoding's $p_e$ (observed or generated) is
floored, as is every $p_a$. An observed peptide with zero support
therefore still gets $p_e =$ \texttt{p\_floor} instead of an exact 0 ---
the numerical floor is the only bound in the system; $\beta$ itself is
unbounded.
\item Zero-encoding rows are handled structurally, not by clamps: a
peptide (observed or drawn) with an empty encoding set has
$\sum \log(\cdot) = 0 \Rightarrow s = 0$ exactly, because the padded
support slots are masked out before the sum.
\item L2 penalty on network weights and $p_o$ (\texttt{reg\_lambda},
default 0); $\beta$ is never penalized.
\item Checkpoint selection: the full objective (exact positive term $+$
importance-sampled unlabeled term) is evaluated every
\texttt{validation\_interval} (10) steps on all observed peptides plus
the fixed validation draws; the best checkpoint is restored at the end.
AdamW (lr 0.01, weight decay 0.01) with ReduceLROnPlateau (patience 200
steps, factor 0.5).
\end{itemize}

\section{Key defaults}\label{sec:defaults}

\begin{table}[htbp]
\centering
\caption{Key configuration defaults}
\label{tab:defaults}
{\footnotesize
\begin{tabular}{@{}>{\raggedright\arraybackslash}p{5.9cm}>{\raggedright\arraybackslash}p{2.4cm}>{\raggedright\arraybackslash}p{6.2cm}@{}}
\toprule
setting & default & meaning \\
\midrule
\texttt{pool.fraction} & 0.01 & pool $=$ first whole time columns
holding $\ge 1\%$ of counts \\
\texttt{observation.obs\_dist\_max} & 2 & batch sequences farther from
pool are excluded \\
\texttt{observation.support\_max\_dist} & 3 & parents beyond this
distance contribute no support \\
\texttt{proposal.n\_max} & 3 & max mutations per draw / encoding
nt-change budget \\
\texttt{proposal.lam} & auto & $\lambda = L_{\text{nuc}} \cdot
\mu_{\text{total}}$ \\
\texttt{mutation\_model} & \texttt{ti90\_tv5},
$\mu_{\text{total}} = 2\times10^{-4}$ & \texttt{sars\_cov2} spectrum
opt-in \\
\texttt{training.sgd\_batch\_size} & 128 & accepted draws per step \\
\texttt{training.fixed\_validation\_size} & 500 & fixed validation
draws \\
\texttt{training.early\_stop\_checks} & 100 & $\times$ 10 steps $=$
1000-step patience \\
\texttt{training.p\_floor} & $10^{-8}$ & clamp floor for $p_e, p_o,
p_a$ \\
\bottomrule
\end{tabular}
}
\end{table}

\clearpage
\section{Differences from the June reference at a glance}\label{sec:diffs}

\begin{table}[htbp]
\centering
\caption{June one-batch version vs.\ the current pool + batch pipeline}
\label{tab:diffs}
{\footnotesize
\begin{tabular}{@{}>{\raggedright\arraybackslash}p{3.0cm}>{\raggedright\arraybackslash}p{4.4cm}>{\raggedright\arraybackslash}p{6.9cm}@{}}
\toprule
 & June (one batch) & current (pool + batch) \\
\midrule
parents of support & all observed sequences & pool sequences only \\
proposal space & 1 AA mutation & 1..3 AA mutations, truncated Poisson
count \\
$q_{\text{raw}}$ & $\sum_{p:\,d_{aa}=1} r(p)/19L$ &
$\sum_{m=1}^{3}\sum_{p:\,d_{aa}=m} r(p)\, P_p(n{=}m) \newline
/\, \big(\binom{L_p}{m}\, 19^{m}\big)$ \\
$Z$ & exact, 1-mutation counts & exact, 1..3-mutation counts (MC
fallback) \\
encodings of $x'$ & 1-nt codon changes, one differing residue & minimal
codon paths over $\le 3$ residues, $\le 3$ nt changes total,
codon-context aware \\
support per encoding & transition/transversion rate $\times$ observed
count & exact minimal-pathway rate product over pool parents at
distance 1..3, $\times$ $\text{freq}_0$ \\
self path & --- & explicit self-survival for exact pool members (always
on) \\
observed-side $\hat{\mathcal{Y}}(x)$ & observed encodings only & $+$
extra generated encodings (\texttt{POS\_L1}) \\
surveillance counts in loss & per-count weighting & one term per
peptide \\
\bottomrule
\end{tabular}
}
\end{table}

\end{document}
